\begin{figure}[ht]
  \centering
  \begin{tikzpicture}
    \begin{groupplot}[group style={group size=2 by 1, horizontal sep=1.4cm, y descriptions at=edge left},
        tesi, width=.5\textwidth, height=6.2cm, xmode=log, log basis x=2,
        xtick={1,2,4,8,16,32}, xticklabels={1,2,4,8,16,32}, xmin=1, xmax=32,
        ymin=0, ymax=7000, xlabel={client in parallelo}]
      \nextgroupplot[title={\footnotesize Mac (M2)}, ylabel={ricerche al secondo},
          legend to name=leg:carico, legend columns=3]
        \addplot[color=black, mark=square*, mark size=1.6pt] table[x=client, y=mac_es] {dati/carico.dat};
        \addplot[color=black, mark=triangle*, mark size=2pt, dashed] table[x=client, y=mac_kc] {dati/carico.dat};
        \addplot[color=black, mark=*, mark size=1.6pt, densely dotted] table[x=client, y=mac_kn] {dati/carico.dat};
        \legend{Elasticsearch come l'app, Koskidex nel container, Koskidex nativo}
      \nextgroupplot[title={\footnotesize Secondo computer (Windows, WSL2)}]
        \addplot[color=black, mark=square*, mark size=1.6pt] table[x=client, y=fisso_es] {dati/carico.dat};
        \addplot[color=black, mark=triangle*, mark size=2pt, dashed] table[x=client, y=fisso_kc] {dati/carico.dat};
        \addplot[color=black, mark=*, mark size=1.6pt, densely dotted] table[x=client, y=fisso_kn] {dati/carico.dat};
    \end{groupplot}
  \end{tikzpicture}

  \ref{leg:carico}
  \caption{Ricerche al secondo al crescere dei client che le mandano senza
  pause, sui 10.018 atti dell'albo. Elasticsearch, come lo usa l'app, e Koskidex
  nel container hanno le stesse quattro CPU; Koskidex nativo gira fuori da
  Docker. Sul Mac Koskidex è dopo le correzioni delle allocazioni. Dati:
  \texttt{latex/dati/carico.dat}.}
  \label{fig:carico}
\end{figure}
